\chapter{2002年粤豫苏桂卷}
\section{单选题}

\begin{problem}
函数 \(f(x) = \dfrac{\sin 2x}{\cos x}\) 的最小正周期是（\(\quad\)）

\choices
    {\(\dfrac{\pi}{2}\)}
    {\(\pi\)}
    {\(2\pi\)}
    {\(4\pi\)}
\end{problem}

\begin{problem}
圆 \((x - 1)^2 + y^2 = 1\) 的圆心到直线 \(y = \dfrac{\sqrt{3}}{3}x\) 的距离是（\(\quad\)）

\choices
    {\(\dfrac{1}{2}\)}
    {\(\dfrac{\sqrt{3}}{2}\)}
    {\(1\)}
    {\(\sqrt{3}\)}
\end{problem}

\begin{problem}
不等式 \((1 + x)(1 - |x|) > 0\) 的解集是（\(\quad\)）

\choices
    {\(\{x \mid 0 \leqslant x < 1\}\)}
    {\(\{x \mid x < 0 \text{ 且 } x \neq -1\}\)}
    {\(\{x \mid -1 < x < 1\}\)}
    {\(\{x \mid x < 1 \text{ 且 } x \neq -1\}\)}
\end{problem}

\begin{problem}
在 \((0,2\pi)\) 内，使 \(\sin x > \cos x\) 成立的 \(x\) 的取值范围是（\(\quad\)）

\choices
    {\(\left(\dfrac{\pi}{4},\dfrac{\pi}{2}\right)\cup \left(\pi,\dfrac{5\pi}{4}\right)\)}
    {\(\left(\dfrac{\pi}{4},\pi\right)\)}
    {\(\left(\dfrac{\pi}{4},\dfrac{5\pi}{4}\right)\)}
    {\(\left(\dfrac{\pi}{4},\pi\right) \cup \left(\dfrac{5\pi}{4}, \dfrac{3\pi}{2}\right)\)}
\end{problem}

\begin{problem}
设集合 \(M = \left\{x \mid x = \dfrac{k}{2} + \dfrac{1}{4}, k \in \Z\right\}\)，\(N = \left\{x \mid x = \dfrac{k}{4} + \dfrac{1}{2}, k \in \Z\right\}\)，则（\(\quad\)）

\choices
    {\(M = N\)}
    {\(M \subseteq N\)}
    {\(M \supseteq N\)}
    {\(M \cap N = \varnothing\)}
\end{problem}

\begin{problem}
一个圆锥和一个半球有公共底面，如果圆锥的体积恰好与半球的体积相等，那么这个圆锥轴截面顶角余弦值是（\(\quad\)）

\choices
    {\(\dfrac{3}{4}\)}
    {\(\dfrac{4}{5}\)}
    {\(\dfrac{3}{5}\)}
    {\(-\dfrac{3}{5}\)}
\end{problem}

\begin{problem}
函数 \(f(x) = x|x + a| + b\) 是奇函数的充要条件是（\(\quad\)）

\choices
    {\(ab = 0\)}
    {\(a + b = 0\)}
    {\(a = b\)}
    {\(a^2 + b^2 = 0\)}
\end{problem}

\begin{problem}
已知 \(0 < x < y < a < 1\)，则有（\(\quad\)）

\choices
    {\(\log_a(xy) < 0\)}
    {\(0 < \log_a(xy) < 1\)}
    {\(1 < \log_a(xy) < 2\)}
    {\(\log_a(xy) > 2\)}
\end{problem}

\begin{problem}
函数 \(y = 1 - \dfrac{1}{x - 1}\)（\(\quad\)）

\choices
    {在 \((-1, +\infty)\) 内单调递增}
    {在 \((-1, +\infty)\) 内单调递减}
    {在 \((1, +\infty)\) 内单调递增}
    {在 \((1, +\infty)\) 内单调递减}
\end{problem}

\begin{problem}
极坐标方程 \(\rho = \cos \theta\) 与 \(\rho \cos \theta = \dfrac{1}{2}\) 的图形是（\(\quad\)）

\choices
    {\choicebitmap[width=0.15\paperwidth]{img/2002/yue_yu_su_gui/q10_optA_nolabel.png}}
    {\choicebitmap[width=0.15\paperwidth]{img/2002/yue_yu_su_gui/q10_optB_nolabel.png}}
    {\choicebitmap[width=0.15\paperwidth]{img/2002/yue_yu_su_gui/q10_optC_nolabel.png}}
    {\choicebitmap[width=0.15\paperwidth]{img/2002/yue_yu_su_gui/q10_optD_nolabel.png}}
\end{problem}

\begin{problem}
从正方体的 \(6\) 个面中选取 \(3\) 个面，其中有 \(2\) 个面不相邻的选法共有（\(\quad\)）

\choices
    {\(8\) 种}
    {\(12\) 种}
    {\(16\) 种}
    {\(20\) 种}
\end{problem}

\begin{problem}
据 \(2002\text{年}\) \(3\text{月}\) \(5\text{日}\)九届人大五次会议《政府工作报告》： “\(2001\text{年}\)国内生产总值达到 \(95933\text{亿元}\)，比上年增长 \(7.3\%\)”，如果“十·五”期间（\(2001\text{年}\)－\(2005\text{年}\)）每年的国内生产总值都按此年增长率增长，那么到“十·五”末我国国内年生产总值约为（\(\quad\)）

\choices
    {\(115000\) 亿元}
    {\(120000\) 亿元}
    {\(127000\) 亿元}
    {\(135000\) 亿元}
\end{problem}

\section{填空题}

\begin{problem}
椭圆 \(5x^{2} + ky^{2} = 5\) 的一个焦点是 \((0,2)\)，那么 \(k =\)\fillinblank{}.
\end{problem}

\begin{problem}
\((x^{2} + 1)(x - 2)^{7}\) 展开式中 \(x^{3}\) 的系数是\fillinblank{}.
\end{problem}

\begin{problem}
已知 \(\sin \alpha = \cos 2\alpha\)，\(\alpha \in \left(\dfrac{\pi}{2}, \pi\right)\)，则 \(\tan \alpha =\)\fillinblank{}.
\end{problem}

\begin{problem}
已知 \(f(x) = \dfrac{x^{2}}{1 + x^{2}}\)，那么 \(f(1) + f(2) + f\left(\dfrac{1}{2}\right) + f(3) + f\left(\dfrac{1}{3}\right) + f(4) + f\left(\dfrac{1}{4}\right) =\)\fillinblank{}.
\end{problem}

\section{解答题}

\begin{problem}
已知复数 \(z = 1 + \i\)，求实数 \(a\)，\(b\) 使 \(az + 2b\bar{z} = (a + 2z)^{2}\).
\end{problem}

\begin{problem}
设 \(\{a_{n}\}\) 为等差数列，\(\{b_{n}\}\) 为等比数列，\(a_{1} = b_{1} = 1\)，\(a_{2} + a_{4} = b_{3}\)，\(b_{2}b_{4} = a_{3}\)，分别求出 \(\{a_{n}\}\) 及 \(\{b_{n}\}\) 的前 \(10\) 项的和 \(S_{10}\) 及 \(T_{10}\).
\end{problem}

\begin{problem}
四棱锥 \(P - ABCD\) 的底面是边长为 \(a\) 的正方形，\(PB \perp\) 平面 \(ABCD\).
\begin{enumerate}
\item 若面 \(PAD\) 与面 \(ABCD\) 所成的二面角为 \(60^{\circ}\)，求这个四棱锥的体积；
\item 证明无论四棱锥的高怎样变化，面 \(PAD\) 与面 \(PCD\) 所成的二面角恒大于 \(90^{\circ}\).
\end{enumerate}
\end{problem}

\begin{problem}
设 \(A\)，\(B\) 是双曲线 \(x^{2} - \dfrac{y^{2}}{2} = 1\) 上的两点，点 \(N(1,2)\) 是线段 \(AB\) 的中点.

\begin{enumerate}
\item 求直线 \(AB\) 的方程；
\item 如果线段 \(AB\) 的垂直平分线与双曲线相交于 \(C\)，\(D\) 两点，那么 \(A\)，\(B\)，\(C\)，\(D\) 四点是否共圆？ 为什么？
\end{enumerate}
\end{problem}

\begin{problem}
\begin{enumerate}
\item 给出两块相同的正三角形纸片（如图 1，图 2）. 要求用其中一块剪拼成一个三棱锥模型，另一块剪拼成一个正三棱柱模型，使它们的全面积都与原三角形的面积相等，请设计一种剪拼方法，分别用虚线标示在图 1，图 2 中，并作简要说明；
\item 试比较你剪拼的正三棱锥与正三棱柱的体积的大小；
\item 如果给出的是一块任意三角形的纸片（如图 3），要求剪拼成一个直三棱柱，使它的全面积与给出的三角形的面积相等. 请设计一种剪拼方法，用虚线标示在图 3 中，并作简要说明.

\bitmapfigure[max width=0.74\linewidth]{img/2002/yue_yu_su_gui/q21_fig1.png}
\end{enumerate}
\end{problem}

\begin{problem}
\begin{enumerate}
\item 当 \(b > 0\) 时，若对任意 \(x \in \R\) 都有 \(f(x) \leqslant 1\)，证明：\(a \leqslant 2\sqrt{b}\)；
\item 当 \(b > 1\) 时，证明：对任意 \(x \in [0,1]\)，\(|f(x)| \leqslant 1\) 的充要条件是 \(b - 1 \leqslant a \leqslant 2\sqrt{b}\)；
\item 当 \(0 < b \leqslant 1\) 时，讨论：对任意 \(x \in [0,1]\)，\(|f(x)| \leqslant 1\) 的充要条件.
\end{enumerate}
\end{problem}
